\chapter{Deep learning}
\label{cap:deep}

\begin{inbreve}
Il deep learning (DL) usa reti con molti strati che apprendono rappresentazioni dei dati su più livelli di
astrazione. Il capitolo discute perché la profondità può servire (risultati di \emph{no-flattening}, l'esempio della
parità, il bias induttivo composizionale), i concetti di representation learning (pre-training, autoencoder,
transfer learning) e di rappresentazione distribuita, alcuni risultati recenti e le tecniche che rendono addestrabili
le reti profonde (ReLU, gradient clipping, batch normalization, dropout).
\end{inbreve}

\section{Che cos'è il deep learning}

\begin{definizione}[Deep learning]
Il \termine{deep learning} permette a modelli computazionali composti da più strati di elaborazione di
\textbf{apprendere rappresentazioni dei dati con più livelli di astrazione}.
\end{definizione}

Il DL è un \textbf{quadro generale} che include reti profonde (MLP con molti strati), reti convoluzionali, modelli
generativi (\ingl{deep belief network}), reti ricorrenti e ricorsive, e si contrappone ai modelli \textbf{shallow}
(superficiali, ad esempio con meno di 3 strati). Gli aspetti comuni sono:
\begin{itemize}
  \item \textbf{più strati} di unità di elaborazione \textbf{non lineari};
  \item apprendimento (supervisionato o non supervisionato) di \textbf{rappresentazioni delle feature} in ogni strato,
        con gli strati che formano una \textbf{gerarchia} da feature di basso livello a feature di alto livello;
  \item quindi, una rappresentazione \textbf{gerarchica, sparsa e distribuita} dei dati.
\end{itemize}
Il suo successo pratico dipende da tre fattori: \textbf{grandi dataset}, \textbf{potenza di calcolo} (le GPU) e
\textbf{nuove idee} per facilitare l'addestramento.

\subsection{Astrazione gerarchica}

Come nello XOR (\rif{sec:xor}), lo strato nascosto crea una ri-rappresentazione che semplifica il compito
dello strato successivo; il DL estende questa composizione su molti livelli.
\begin{itemize}
  \item Si passa dal \textbf{feature engineering} (feature estratte a mano, per esperienza) al \textbf{feature
        learning}: si \textbf{apprendono rappresentazioni} dei dati, anche in modo non supervisionato o
        semi-supervisionato.
  \item I livelli di astrazione crescono con gli strati: un'immagine si può rappresentare come vettore di intensità
        dei pixel, oppure come insieme di bordi, regioni di forma particolare, ecc. \textbf{I bordi formano motivi, i
        motivi si assemblano in parti, le parti formano oggetti} (\figrif{fig:gerarchia}).
  \item Funziona al meglio quando l'input ha una \textbf{struttura} (spaziale, temporale, \dots).
\end{itemize}

\begin{figure}[ht]
\centering
\begin{tikzpicture}[node distance=0.55cm]
\node[blocco, minimum width=2.1cm] (a) {pixel};
\node[blocco, minimum width=2.1cm, right=of a] (b) {bordi};
\node[blocco, minimum width=2.1cm, right=of b] (c) {angoli\\e contorni};
\node[blocco, minimum width=2.1cm, right=of c] (d) {parti\\di oggetti};
\node[blocco, minimum width=2.1cm, right=of d, line width=1pt] (e) {identità\\dell'oggetto};
\foreach \u/\v in {a/b,b/c,c/d,d/e} \draw[freccia] (\u) -- (\v);
\node[font=\scriptsize, grigio, below=2pt of a] {input};
\node[font=\scriptsize, grigio, below=2pt of b] {strato 1};
\node[font=\scriptsize, grigio, below=2pt of c] {strato 2};
\node[font=\scriptsize, grigio, below=2pt of d] {strato 3};
\node[font=\scriptsize, grigio, below=2pt of e] {uscita};
\end{tikzpicture}
\caption{Una mappatura difficile (dai pixel all'identità dell'oggetto) scomposta in una serie di mappature semplici
annidate, ciascuna descritta da uno strato.}
\label{fig:gerarchia}
\end{figure}

Il DL risolve una mappatura complicata \textbf{scomponendola in mappature semplici annidate}, e la nuova
rappresentazione semplifica la classificazione nell'ultimo strato. Le feature si \textbf{apprendono dagli esempi}:
con la backpropagation l'errore in uscita si propaga attraverso gli strati e le unità si specializzano nelle feature
utili al task. I primi filtri di una rete per immagini estraggono tipicamente linee orizzontali o verticali, come un
semplice filtro convolutivo.

\subsection{Combinare per generalizzare}

Le feature astratte si \textbf{riusano} ai livelli superiori e si \textbf{combinano} per generalizzare a casi mai
visti: se la rete ha visto donne senza occhiali e uomini con occhiali, può rappresentare anche una donna con gli
occhiali. Operazioni sulle rappresentazioni vettoriali apprese lo mostrano: ``uomo con occhiali $-$ uomo $+$ donna
$=$ donna con occhiali''. Lo stesso accade nel linguaggio:
\begin{itemize}
  \item $\mathrm{Rep}(\text{re}) - \mathrm{Rep}(\text{maschio}) + \mathrm{Rep}(\text{femmina})\approx\mathrm{Rep}(\text{regina})$:
        la differenza tra ``re'' e ``maschio'' cattura il concetto di monarchia;
  \item $\mathrm{Rep}(\text{Parigi}) - \mathrm{Rep}(\text{Francia}) + \mathrm{Rep}(\text{Polonia})\approx\mathrm{Rep}(\text{Varsavia})$:
        la differenza cattura il concetto di capitale.
\end{itemize}

\section{Servono davvero molti strati?}
\label{sec:no-flattening}

\begin{idea}[No-flattening]
\emph{Una funzione che si rappresenta in modo compatto con un'architettura profonda può richiedere un'architettura
molto grande per essere rappresentata con una insufficientemente profonda.}
\end{idea}

Un circuito logico a \textbf{due strati} rappresenta \textbf{qualsiasi} funzione booleana in forma normale disgiuntiva
(porte AND nel primo strato, eventualmente con input negati, e una porta OR nel secondo), ma con profondità due
\textbf{la maggior parte delle funzioni booleane richiede un numero esponenziale di porte}.

\begin{esempio}[La funzione di parità]
La \textbf{parità dispari} estende lo XOR a $N$ input: vale 1 se e solo se il numero di 1 tra gli $N$ bit è dispari
(si assumono disponibili anche i NOT).
\begin{itemize}
  \item \textbf{Con 2 strati} bisogna elencare tutte le configurazioni positive. Per 3 input:
        $x_1x_2x_3 + x_1\bar x_2\bar x_3 + \bar x_1x_2\bar x_3 + \bar x_1\bar x_2x_3$ (1 se l'input è 111, 100, 010 o
        001), cioè 4 AND e 1 OR, 5 porte. In generale le configurazioni con un numero dispari di 1 sono metà delle
        $2^N$, quindi servono $2^N/2$ porte AND più l'OR finale: $2^{N-1}+1$ porte, cioè 9 per $N=4$ e
        \textbf{129} per $N=8$, un numero \textbf{esponenziale}.
  \item \textbf{Con $\log N$ strati} si usa un \textbf{albero binario completo} di XOR (\figrif{fig:parita}): un
        albero con $N$ foglie ha $N-1$ nodi interni, ognuno uno XOR da 3 porte (AND e OR come nel
        \rif{sec:xor}), per un totale di $3(N-1)$ porte, un numero \textbf{polinomiale}: 21 porte per $N=8$
        contro 129.
\end{itemize}
\end{esempio}

\begin{figure}[ht]
\centering
\begin{tikzpicture}[x=0.95cm, y=1.05cm, xr/.style={circle, draw, fill=grigiochiaro, inner sep=1pt, font=\scriptsize, minimum size=6mm}]
\foreach \i in {1,...,8} \node[font=\small] (x\i) at (\i,0) {$x_\i$};
\foreach \k in {1,...,4} { \pgfmathtruncatemacro{\a}{2*\k-1} \pgfmathtruncatemacro{\b}{2*\k}
  \node[xr] (p\k) at (2*\k-0.5,1) {XOR}; \draw (x\a) -- (p\k); \draw (x\b) -- (p\k); }
\node[xr] (q1) at (2.5,2) {XOR}; \draw (p1) -- (q1); \draw (p2) -- (q1);
\node[xr] (q2) at (6.5,2) {XOR}; \draw (p3) -- (q2); \draw (p4) -- (q2);
\node[xr] (r) at (4.5,3) {XOR}; \draw (q1) -- (r); \draw (q2) -- (r);
\draw[->] (r) -- (4.5,3.7) node[above, font=\small] {parità};
\node[font=\footnotesize, align=left, anchor=west] at (9,1.5) {$\log_2 8 = 3$ livelli di XOR,\\$8-1 = 7$ XOR da 3 porte:\\$21$ porte in tutto};
\end{tikzpicture}
\caption{Parità di 8 bit con un albero di XOR: un numero di porte lineare in $N$, contro le $2^{N-1}+1$ di un
circuito a due strati.}
\label{fig:parita}
\end{figure}

Questo vantaggio \textbf{non vale per tutte le classi di funzioni}: caratterizzare quelle che traggono un vantaggio
(esponenziale) dalla profondità è un problema di ricerca aperto. Per le reti neurali i risultati sono di due tipi,
entrambi \textbf{sull'efficienza} e non sull'espressività (per l'approssimazione universale basta uno strato):
\begin{enumerate}
  \item \textbf{reti shallow}: esistono casi in cui un solo strato nascosto richiede un numero \textbf{esponenziale} di
        unità (rispetto alla dimensione $n$ dell'input) o di pesi non nulli. Nel caso binario le funzioni su
        $\{0,1\}^n$ sono $2^{2^n}$, e nel caso peggiore può servire un'unità per ogni configurazione di input da
        distinguere, come nella parità. Non è solo un problema di dimensione: significa anche che diventa
        \textbf{difficile apprendere con pochi esempi}, ed è questo il punto più rilevante per il ML;
  \item \textbf{no-flattening delle reti profonde}: esistono famiglie di funzioni approssimabili in modo efficiente con
        profondità maggiore di $d$ che richiedono un modello molto più grande (fino a un fattore esponenziale) se la
        profondità è al più $d$.
\end{enumerate}
Non c'è però garanzia che un dato task abbia questa proprietà.

\subsection{Composizionalità e bias induttivo}
\label{sec:bias-deep}

Le reti profonde sfruttano la \textbf{composizionalità delle rappresentazioni interne}: i concetti semplici di uno
strato sono usati come \textbf{primitive} dallo strato successivo per rappresentare concetti più complessi, evitando la
rappresentazione (e l'apprendimento) esplicita e combinatoria di tutte le combinazioni. Si imparano le singole parole
o i sotto-pattern, non direttamente tutte le loro combinazioni.

\begin{idea}[Profondo non significa grande]
Le reti profonde non vanno viste come modelli \emph{enormi}: sono modelli \textbf{compatti} rispetto a modelli
superficiali potenzialmente molto più grandi (anche esponenzialmente) per lo stesso task. Meno unità e meno pesi
aiutano anche l'apprendimento: permettono di generalizzare bene con \textbf{meno esempi}. Una piccola rete profonda
può sostituire un'enorme rete superficiale.
\end{idea}

Scegliere un modello profondo codifica una convinzione molto generale: che la funzione da apprendere sia una
\textbf{composizione di funzioni più semplici} (o di fattori di variazione annidati). Se il task corrisponde a questo
bias, la generalizzazione migliora proprio grazie ai molti strati: a parità di numero di parametri, reti più profonde
generalizzano meglio, mentre aumentare i parametri senza aumentare la profondità non aiuta. Task tipicamente
gerarchici sono le \textbf{immagini} (composizione di parti) e il \textbf{linguaggio}, ma non tutti i dati e i task. È
comunque un bias molto più \textbf{generale} di vincoli ad hoc (le $\phi$ scelte a mano di una LBE, le feature
ingegnerizzate); nulla però funziona senza assunzioni.

\subsection{Maledizione della dimensionalità e bias}

\begin{itemize}
  \item Molti learner si basano su \textbf{approssimazioni locali} (un prior di smoothness locale): K-NN, kernel
        locali, alberi di decisione. Per distinguere $O(k)$ regioni richiedono $O(k)$ esempi, e in alta dimensione
        servono moltissimi esempi per coprire lo spazio (\rif{sec:curse}).
  \item Per distinguere un numero \textbf{esponenziale} di regioni (in $k$) con $O(k)$ esempi, cioè per generalizzare
        \textbf{non localmente}, servono ipotesi aggiuntive sulla distribuzione dei dati.
  \item Il DL sceglie un bias \textbf{generale}, legato alla composizione di funzioni: i dati sono generati dalla
        composizione di fattori, su più livelli. Quando l'ipotesi è vera si ottiene un guadagno potenzialmente
        esponenziale tra numero di esempi e numero di regioni distinguibili.
  \item La distribuzione delle immagini e dei testi reali sembra concentrata su una varietà (\ingl{manifold}) che
        occupa una piccola frazione del volume di tutte le immagini o i testi possibili.
\end{itemize}

Quanti strati e quante unità? Non c'è una risposta generale: è un problema di model selection. Le reti profonde usano
spesso meno unità per strato, quindi meno parametri e meno dati per generalizzare, ma molti strati possono essere
\textbf{più difficili da ottimizzare} (\rif{sec:tecniche-deep}); con modelli così grandi l'esperienza conta molto nel
fissare gli aspetti architetturali.

\section{Representation learning}
\label{sec:representation}

\begin{definizione}[Representation learning]
Il \termine{representation learning} è l'insieme dei metodi che permettono a una macchina di ricevere dati grezzi e
scoprire automaticamente le rappresentazioni necessarie per il rilevamento o la classificazione.
\end{definizione}

I metodi di DL sono metodi di representation learning con più livelli di rappresentazione, ma il concetto vale per
molti modelli (le reti neurali in generale) e riguarda soprattutto i \textbf{dati grezzi} (immagini, testo).
\begin{itemize}
  \item Molti compiti sono facili o difficilissimi \textbf{a seconda di come l'informazione è rappresentata}.
  \item Nel ML \textbf{una buona rappresentazione è quella che rende più facile il compito di apprendimento
        successivo}.
  \item Progettare le feature a mano è difficile ed equivale a scegliere a mano le $\phi(\x)$ di una LBE;
        l'apprendimento supervisionato di un MLP produce invece una rappresentazione $\phi(\x,\w)$ in ogni strato
        nascosto, appresa trovando il miglior $\w$.
\end{itemize}
Oltre all'addestramento completo di MLP e CNN, le rappresentazioni nascoste si possono \textbf{ottenere} con
l'apprendimento non supervisionato (pre-training) e \textbf{sfruttare} per altri task (transfer learning).

\subsection{Pre-training e autoencoder}

Il \termine{pre-training} strato per strato non supervisionato fu il primo approccio che rese possibile addestrare
reti profonde (a parte CNN e RNN): si cattura la distribuzione dell'input strato per strato e si inizializza così la
rete, invece di addestrarla subito \ingl{end-to-end} con la discesa del gradiente.

\begin{definizione}[Autoencoder]
Un \termine{autoencoder} è una rete addestrata a \textbf{copiare il proprio input in uscita}. Ha uno strato nascosto
$\h$ che descrive un \textbf{codice} dell'input, ed è composto da un \textbf{encoder} $\h = f(\x)$ e da un
\textbf{decoder} che produce la ricostruzione $\vv r = g(\h)$. È apprendimento non supervisionato.
\end{definizione}

\begin{figure}[ht]
\centering
\begin{tikzpicture}[x=1.3cm, y=0.55cm]
\foreach \i in {1,...,6} \node[neur, minimum size=5mm] (a\i) at (0,3.5-\i) {};
\foreach \i in {1,...,3} \node[neurp, minimum size=5mm] (h\i) at (2,2-\i) {};
\foreach \i in {1,...,6} \node[neur, minimum size=5mm] (r\i) at (4,3.5-\i) {};
\foreach \i in {1,...,6} \foreach \j in {1,...,3} { \draw[grigio] (a\i) -- (h\j); \draw[grigio] (h\j) -- (r\i); }
\node[font=\small] at (0,3.5) {$\x$}; \node[font=\small] at (2,1.6) {$\h$}; \node[font=\small] at (4,3.5) {$\vv r\approx\x$};
\node[font=\footnotesize] at (1,-4.4) {encoder $f$}; \node[font=\footnotesize] at (3,-4.4) {decoder $g$};
\end{tikzpicture}
\caption{Un autoencoder \emph{undercomplete}: lo strato nascosto, più piccolo dell'input, è costretto a catturare le
feature più salienti per poter ricostruire l'input.}
\label{fig:autoencoder}
\end{figure}

\begin{itemize}
  \item \textbf{Undercomplete}: strato nascosto \textbf{più piccolo} dell'input; il vincolo architetturale forza la rete
        a catturare le feature più salienti (con decoder lineare e loss quadratica si ottiene la PCA).
  \item \textbf{Overcomplete}: strato nascosto \textbf{più grande} dell'input, con una regolarizzazione che impone
        sparsità, robustezza al rumore o altre proprietà, oltre alla banale capacità di copiare.
  \item \textbf{Denoising autoencoder}: si addestra a ricostruire l'input pulito da una versione corrotta (una buona
        rappresentazione si ottiene in modo robusto da un input corrotto ed è utile per recuperare l'input pulito).
\end{itemize}

\begin{algoritmo}[Pre-training strato per strato]
\begin{enumerate}
  \item Addestrare il primo strato come \textbf{autoassociatore}, minimizzando l'errore di ricostruzione dell'input
        grezzo (bastano esempi non etichettati).
  \item Usare le uscite delle unità nascoste come input di un altro strato, anch'esso addestrato come autoassociatore.
  \item Ripetere fino al numero di strati desiderato.
  \item Usare l'uscita dell'ultimo strato nascosto come input di uno strato supervisionato.
  \item \textbf{Fine-tuning} di tutti i parametri rispetto al criterio supervisionato.
\end{enumerate}
\end{algoritmo}

Il pre-training funziona sia come \textbf{inizializzazione} (si spera di portare i parametri in una regione da cui la
discesa locale raggiunge un buon ottimo) sia come \textbf{regolarizzazione}, non nel senso di un modello più semplice
ma di feature che catturano le regolarità della distribuzione degli input; riduce anche la \textbf{varianza} della
stima. Oggi è largamente abbandonato per l'addestramento da zero, perché le tecniche moderne bastano per un buon
addestramento end-to-end, ma resta utile dove permette di sfruttare grandi quantità di dati non etichettati (ad esempio
l'NLP). Architetture collegate: gli \textbf{autoencoder variazionali} (lo spazio latente è una distribuzione), le
architetture \textbf{encoder-decoder} per la traduzione, i \textbf{transformer} (\rif{sec:transformer}), usati con il
solo encoder (BERT) o il solo decoder per la generazione autoregressiva (GPT).

\subsection{Transfer learning}
\label{sec:transfer}

Si usa la rappresentazione appresa da un modello per \textbf{migliorarne un altro}, assumendo che esistano feature utili
in contesti o task diversi. Forme diverse: dall'autoencoder al classificatore; \textbf{multi-task learning} (stessi
input, target diversi: una regolarizzazione implicita); \textbf{domain adaptation} (domini di input diversi con feature
condivise); riuso di \textbf{modelli pre-addestrati} su dataset più grandi.

\begin{esempio}[Riusare una CNN pre-addestrata]
Una CNN addestrata su circa 1,2 milioni di immagini in 1000 categorie ha appreso rappresentazioni ricche, con primitive
di basso livello condivise tra immagini diverse. Per un nuovo task si \textbf{sostituiscono gli ultimi strati} e si
addestrano solo quelli con poche immagini nuove, oppure si fa il \textbf{fine-tuning} di tutta la rete.
\end{esempio}

\section{Rappresentazioni distribuite}

\begin{definizione}[Rappresentazione distribuita]
In una \termine{rappresentazione distribuita} gli elementi (le feature) non sono mutuamente esclusivi, e le loro molte
configurazioni corrispondono alle variazioni osservate nei dati.
\end{definizione}

\begin{table}[ht]
\centering\small
\caption{Rappresentazione simbolica (one-hot) e distribuita.}
\begin{tabularx}{0.95\linewidth}{XX}
\toprule
Simbolica (one-hot) & Distribuita \\
\midrule
una sola unità vale 1, le altre 0 & vettori di numeri reali (densi) \\
distanza sempre $\sqrt2$ tra concetti diversi & ogni concetto è l'insieme delle attivazioni di tutte le unità \\
 & la similarità si può condividere tra più di due concetti; si condividono feature apprese; la distanza riflette il significato \\
\bottomrule
\end{tabularx}
\end{table}

Qui non si classifica, si \textbf{rappresenta}. L'apprendimento agisce sulla rappresentazione \textbf{interna}: una
rappresentazione distribuita in \textbf{input} si usa solo se la conoscenza di dominio permette di fissarla (fissarla
arbitrariamente può danneggiare il task); di solito si fornisce un input one-hot e si lascia che il modello sviluppi
internamente la rappresentazione distribuita.

\subsection{Contare la differenza}

Quattro concetti: auto blu, bici blu, auto rossa, bici rossa. Con una rappresentazione one-hot servono 4 unità; con una
distribuita ne bastano \textbf{2} ($2^2$ concetti), una per il colore e una per il tipo, e l'unità del ``rosso'' impara
il rosso da immagini sia di auto sia di bici. In generale, con $n$ feature a $k$ valori una rappresentazione
distribuita descrive $k^n$ concetti, contro gli $n$ di una one-hot: si può assegnare un codice a un numero
\textbf{esponenziale} di regioni. Sono invece non distribuiti il K-NN (uno o pochi prototipi per input), i kernel RBF
stretti (che si riducono al K-NN), le miscele di gaussiane, gli alberi di decisione (si attiva una sola foglia).

\begin{idea}[Attributi condivisi]
Separando (\ingl{disentangling}) i concetti di colore e di tipo, si impara la distinzione auto/bici o il colore
\textbf{senza vedere esempi di tutte le combinazioni}: avendo visto gli altri tre casi, si tratta facilmente la bici
blu, che per una rappresentazione one-hot sarebbe del tutto nuova. Con le parole, una rappresentazione in cui ``cane''
e ``gatto'' condividono feature permette di trattarle in modo simile (le frasi con ``gatto'' informano quelle con
``cane'') e di imparare automaticamente che sono simili perché compaiono in contesti simili: è il \termine{word
embedding}. Nello spazio one-hot ogni parola è a distanza $\sqrt2$ da tutte le altre; nello spazio appreso parole di
significato simile (nomi di paesi, anni) diventano vicine.
\end{idea}

Il rovescio della medaglia è l'\textbf{interpretabilità}: una rappresentazione distribuita è meno facile da
interpretare di una localista. Il DL sfrutta rappresentazioni distribuite \textbf{attraverso molti strati}: la
composizionalità porta un ulteriore guadagno di efficienza statistica, per \textbf{due vantaggi potenzialmente
esponenziali} rispetto ai modelli non distribuiti e superficiali.

\section{Altri risultati recenti}
\label{sec:double-descent}

\begin{itemize}
  \item \textbf{Smoothness dalla struttura}: a parità di parametri, le reti più profonde impongono più smoothness di
        quelle superficiali (ogni strato lavora sulla superficie già liscia prodotta dal precedente): vincoli
        \textbf{impliciti} anziché espliciti.
  \item \textbf{Il puzzle del non-overfitting}: reti sovra-parametrizzate capaci di azzerare l'errore di training anche
        su etichette \textbf{casuali} generalizzano comunque bene sui dati reali. Una spiegazione è che la discesa del
        gradiente imponga una \textbf{regolarizzazione implicita}, convergendo verso soluzioni a \textbf{margine massimo}.
  \item \termine{Double descent}: al crescere della dimensione del modello (o dei dati, o del tempo di addestramento)
        l'errore di test prima scende, poi sale (la classica U), raggiunge un picco attorno alla \textbf{soglia di
        interpolazione} (dove l'errore di training si azzera) e poi \textbf{scende di nuovo} (\figrif{fig:double-descent}).
        L'effetto è spesso evitato con una regolarizzazione accurata e non è ancora pienamente compreso; non riguarda
        necessariamente i modelli piccoli.
  \item \textbf{Ridondanza}: nei modelli profondi una gran parte dei pesi si può predire dagli altri o eliminare
        (pruning, quantizzazione) senza perdere accuratezza.
  \item \textbf{Lottery ticket hypothesis}: le buone prestazioni dipenderebbero da un'inizializzazione fortunata di una
        o più sotto-reti, e le reti grandi ne contengono esponenzialmente di più; riaddestrando la rete potata con la
        stessa inizializzazione si ottengono prestazioni simili. Il ruolo della casualità è il tema del
        capitolo~\ref{cap:random}.
\end{itemize}

\begin{figure}[ht]
\centering
\begin{tikzpicture}
\begin{axis}[stile, width=8.6cm, height=4.8cm, domain=0.3:10, samples=200, xmin=0, xmax=10.4, ymin=0, ymax=1.15,
  xtick={4}, xticklabels={soglia di interpolazione}, ytick=\empty, xlabel={dimensione del modello}, ylabel={errore},
  legend pos=outer north east]
\addplot[thick, blu] {max(0.02, 0.6*exp(-x/1.4) - 0.04*(x>4))}; \addlegendentry{training}
\addplot[thick, rosso] {0.32 + 0.45*exp(-x/0.9) + 0.55*exp(-(x-4)^2/0.5) - 0.012*x}; \addlegendentry{test}
\draw[densely dashed, grigio] (axis cs:4,0) -- (axis cs:4,1.05);
\node[font=\footnotesize, grigio] at (axis cs:1.8,1.08) {regime classico};
\node[font=\footnotesize, grigio, align=center] at (axis cs:7.4,1.0) {regime\\sovra-parametrizzato};
\end{axis}
\end{tikzpicture}
\caption{Il fenomeno della \emph{double descent} (andamento qualitativo): dopo il picco alla soglia di interpolazione
l'errore di test torna a scendere.}
\label{fig:double-descent}
\end{figure}

\section{Tecniche}
\label{sec:tecniche-deep}

Molti strati possono essere difficili da ottimizzare: l'enfasi è sul miglioramento della discesa del gradiente,
sulla regolarizzazione (le reti sono grandi) e sullo sfruttamento dei dati.

\begin{idea}[Tecniche chiave]
\begin{itemize}
  \item \textbf{SGD con momentum} (mini-batch, decadimento di $\eta$) o \textbf{Adam} (\rif{sec:learning-rate}).
  \item Attivazioni \textbf{ReLU} nelle unità nascoste.
  \item Loss di massima verosimiglianza (\textbf{cross-entropy}) con \textbf{softmax} in uscita (anche per evitare
        saturazione e gradienti piccoli).
  \item Regolarizzazione: early stopping, weight decay, \textbf{dropout}, \textbf{batch normalization}.
\end{itemize}
\end{idea}

\subsection{Vanishing ed exploding gradient}

Retropropagando il gradiente attraverso molti strati, se i pesi sono piccoli il gradiente si \textbf{riduce
esponenzialmente} (\termine{vanishing gradient}), anche per via delle derivate delle sigmoidi; se sono grandi
\textbf{cresce esponenzialmente} (\termine{exploding gradient}). Le attivazioni come la tanh hanno derivate in
$(0,1]$, e la regola della catena moltiplica $d$ di questi fattori per raggiungere gli strati vicini all'input: il
delta decresce esponenzialmente con $d$ e quegli strati si addestrano lentissimamente.

\begin{proprieta}{Schizzo formale}
Se $\h_l = F_l(\h_{l-1})$ e $F = F_M\circ\dots\circ F_1$, allora
\[
\frac{\partial F}{\partial w_l} = \frac{\partial F_M}{\partial\h_{M-1}}\cdots\frac{\partial F_l}{\partial w_l},
\qquad \left\lvert\frac{\partial F}{\partial w_l}\right\rvert\approx\rho^{M-l+1}\ \text{ se }\ \left\lVert\frac{\partial F_i}{\partial\h_i}\right\rVert\approx\rho .
\]
Con $\rho<1$ il gradiente svanisce, con $\rho>1$ esplode.
\end{proprieta}

Gli approcci vanno da R-prop alle connessioni \ingl{shortcut}, alle reti randomizzate, alle ReLU.

\subsubsection{Gradient clipping}
La moltiplicazione ripetuta può creare \textbf{scogliere} (\ingl{cliff}) nella funzione di costo: derivate altissime
che ``catapultano'' i pesi molto lontano. Il \termine{gradient clipping} limita la norma del gradiente mantenendone la
direzione:
\[
\text{se }\norm{\vv g}>v\ \text{ allora }\ \vv g\leftarrow\frac{v\,\vv g}{\norm{\vv g}},
\]
con $v$ soglia sulla norma.

\subsection{ReLU}
\label{sec:relu}

\begin{definizione}[ReLU]
\[
f(x) = \max(0,x) = \begin{cases}0 & x<0\\ x & x\ge0.\end{cases}
\]
\end{definizione}

\begin{itemize}
  \item \textbf{Propagazione del gradiente} migliore attraverso molti strati: niente saturazione, perché $f'(x)=1$ per
        ogni net positivo.
  \item \textbf{Attivazione sparsa}: in una rete inizializzata a caso circa metà delle unità nascoste è attiva.
  \item \textbf{Calcolo efficiente}: solo confronti, somme e prodotti.
\end{itemize}
La ReLU \textbf{non è differenziabile in 0}, ma nelle librerie si assume derivata 0 o 1: è un'approssimazione sicura,
perché è improbabile valutare esattamente $f(0)$. Conviene iniziare con net positivi (ad esempio bias piccoli e
positivi, $0{,}1$), così che le unità siano inizialmente attive. Varianti come \textbf{ELU}
($\alpha(e^x-1)$ per $x\le0$, $x$ per $x>0$) e \textbf{Leaky ReLU} ($0{,}01x$ per $x<0$, $x$ per $x\ge0$) apprendono
anche sulla parte sinistra (le unità non si ``spengono'') e hanno attivazione media più vicina a zero
(\figrif{fig:relu}).

\begin{figure}[ht]
\centering
\begin{tikzpicture}
\begin{axis}[stile, width=7.4cm, height=4.6cm, domain=-3:3, samples=120, xmin=-3, xmax=3, ymin=-1.2, ymax=3.1,
  xtick={-2,0,2}, ytick={-1,0,1,2,3}, xlabel={$x$}, legend pos=outer north east]
\addplot[very thick, inchiostro] {max(0,x)}; \addlegendentry{ReLU}
\addplot[thick, rosso, densely dashed] {x>0 ? x : 0.1*x}; \addlegendentry{Leaky ReLU (pendenza $0{,}1$)}
\addplot[thick, blu] {x>0 ? x : exp(x)-1}; \addlegendentry{ELU ($\alpha=1$)}
\addplot[thick, verde] {ln(1+exp(x))}; \addlegendentry{softplus}
\end{axis}
\end{tikzpicture}
\caption{La ReLU e alcune varianti (la Leaky ReLU è disegnata con pendenza $0{,}1$ invece di $0{,}01$ per renderla
visibile).}
\label{fig:relu}
\end{figure}

\subsection{Batch normalization}

La \termine{batch normalization} normalizza, in ogni strato, le attivazioni del mini-batch a media zero e varianza
unitaria, calcolandone le statistiche sul mini-batch; la trasformazione è inclusa nella backpropagation. Come
normalizzare l'input è pratica standard, così si mantengono normalizzati i dati che fluiscono negli strati intermedi.
Effetti: \textbf{regolarizzazione} (per il rumore introdotto dal mini-batch) e \textbf{apprendimento più rapido}.

\subsection{Dropout}
\label{sec:dropout}

Il \termine{dropout} seleziona casualmente una sotto-rete a ogni passo di addestramento. Si spiega con due concetti
noti: \textbf{ensemble} (un bagging implicito ed economico di sotto-reti) e \textbf{regolarizzazione}. Il bagging
(\rif{sec:ensemble}) media modelli ad alta varianza e basso bias; il dropout sfrutta questo principio con un numero
\textbf{esponenziale} di sotto-reti (\figrif{fig:dropout}), avendo a test time \textbf{una sola rete}, e massimizza la
diversità dell'ensemble evitando co-adattamenti tra le unità.

\begin{figure}[ht]
\centering
\begin{tikzpicture}[x=0.42cm, y=0.42cm, nd/.style={circle, draw, inner sep=0pt, minimum size=2.3mm}, ndx/.style={circle, draw=grigio!50, densely dotted, inner sep=0pt, minimum size=2.3mm}]
\foreach \m [count=\q from 0] in {0,...,15} {
  \pgfmathtruncatemacro{\col}{mod(\q,8)} \pgfmathtruncatemacro{\row}{\q/8}
  \pgfmathtruncatemacro{\ia}{mod(\m,2)} \pgfmathtruncatemacro{\ib}{mod(floor(\m/2),2)}
  \pgfmathtruncatemacro{\ha}{mod(floor(\m/4),2)} \pgfmathtruncatemacro{\hb}{mod(floor(\m/8),2)}
  \begin{scope}[xshift=\col*1.9cm, yshift=-\row*2.5cm]
    \coordinate (A) at (0,0); \coordinate (B) at (2,0); \coordinate (C) at (0,2); \coordinate (D) at (2,2); \coordinate (O) at (1,4);
    \ifnum\ia=0 \node[nd] at (A) {}; \else \node[ndx] at (A) {}; \fi
    \ifnum\ib=0 \node[nd] at (B) {}; \else \node[ndx] at (B) {}; \fi
    \ifnum\ha=0 \node[nd] at (C) {}; \else \node[ndx] at (C) {}; \fi
    \ifnum\hb=0 \node[nd] at (D) {}; \else \node[ndx] at (D) {}; \fi
    \node[nd, fill=grigiochiaro] at (O) {};
    \pgfmathtruncatemacro{\ac}{(\ia==0)&&(\ha==0)} \ifnum\ac=1 \draw (A) -- (C); \fi
    \pgfmathtruncatemacro{\ad}{(\ia==0)&&(\hb==0)} \ifnum\ad=1 \draw (A) -- (D); \fi
    \pgfmathtruncatemacro{\bc}{(\ib==0)&&(\ha==0)} \ifnum\bc=1 \draw (B) -- (C); \fi
    \pgfmathtruncatemacro{\bd}{(\ib==0)&&(\hb==0)} \ifnum\bd=1 \draw (B) -- (D); \fi
    \ifnum\ha=0 \draw (C) -- (O); \fi
    \ifnum\hb=0 \draw (D) -- (O); \fi
  \end{scope}
}
\end{tikzpicture}
\caption{Le 16 sotto-reti di una rete con 2 input, 2 unità nascoste e un'uscita, ottenute rimuovendo sottoinsiemi
delle unità non di uscita. Alcune non collegano più l'input all'uscita, ma in reti grandi è improbabile che non
resti alcun percorso.}
\label{fig:dropout}
\end{figure}

\begin{algoritmo}[Dropout]
\begin{itemize}
  \item Ogni volta che si carica un esempio in un mini-batch si campiona una \textbf{maschera binaria} per tutte le
        unità di input e nascoste, indipendentemente per ogni unità; la probabilità di includere un'unità è un
        iperparametro (tipicamente $0{,}8$ per l'input e $0{,}5$ per le unità nascoste).
  \item Ogni unità è moltiplicata per la sua maschera: zero equivale a rimuoverla, e si ottiene una sotto-rete.
  \item Forward, backpropagation e aggiornamento procedono come al solito sulla sotto-rete.
  \item Le unità rimosse vengono poi reinserite con i loro pesi.
\end{itemize}
\end{algoritmo}

Solo una piccola frazione delle sotto-reti viene addestrata, ciascuna per un solo passo, ma la \textbf{condivisione dei
parametri} porta anche le altre a buoni parametri. A test time si approssima la predizione dell'ensemble con la
\ingl{weight scaling inference rule}: si moltiplicano i pesi per la probabilità di inclusione, così che l'uscita attesa
di ogni unità sia la stessa dell'addestramento. L'effetto regolarizzante: riduce le interazioni tra le unità e le
diversifica; riduce la varianza come il bagging; inserisce \textbf{rumore strutturato} (ad esempio riconoscere un volto
senza l'unità che rileva il naso); soprattutto regolarizza ogni unità a essere una feature buona \textbf{in molti
contesti}. Per la regressione lineare equivale a un weight decay $L^2$ con un $\lambda$ diverso per ogni input.

Altre tecniche: la regolarizzazione $L^1$ (\rif{sec:l1}), che elimina feature; le \termine{Generative Adversarial
Network} (GAN), due reti in competizione, una che \textbf{genera} candidati e una che li \textbf{valuta} distinguendoli
dai dati veri.

\begin{puntichiave}
\begin{itemize}
  \item DL: molti strati non lineari che apprendono una gerarchia di rappresentazioni distribuite.
  \item No-flattening: funzioni compatte con molti strati possono richiedere modelli esponenzialmente grandi con pochi
        strati (parità: $2^{N-1}+1$ porte contro $3(N-1)$).
  \item Bias induttivo: la funzione è una composizione di funzioni più semplici.
  \item Representation learning: pre-training (autoencoder), transfer learning. Rappresentazioni distribuite: $k^n$
        configurazioni, attributi condivisi.
  \item Tecniche: ReLU, cross-entropy con softmax, SGD con momentum/Adam, gradient clipping, batch normalization,
        dropout (ensemble implicito di sotto-reti).
\end{itemize}
\end{puntichiave}

\begin{domande}
\begin{enumerate}
  \item Cosa si intende per deep learning? Quali sono gli aspetti comuni ai vari approcci?
  \item Servono davvero molti strati? Discutere l'esempio della parità e i risultati di no-flattening.
  \item Qual è il bias induttivo dei modelli profondi?
  \item Cos'è il representation learning? Pre-training e transfer learning. Cos'è un autoencoder?
  \item Rappresentazioni simboliche contro distribuite: vantaggi delle seconde.
  \item Cos'è il vanishing gradient? Come si affronta?
  \item Descrivere il dropout e spiegarne l'effetto regolarizzante.
\end{enumerate}
\end{domande}
